\documentclass[12pt,a4paper]{book}

% Encodage et langue
\usepackage[utf8]{inputenc}
\usepackage[T1]{fontenc}
\usepackage[french]{babel}

% Mise en page
\usepackage{geometry}
\geometry{margin=2.5cm}

% Tableaux et listes
\usepackage{tabularx}
\usepackage{enumitem}

% Maths et symboles
\usepackage{amsmath,amssymb}

% Liens
\usepackage{hyperref}
\hypersetup{
    colorlinks=true,
    linkcolor=blue,
    urlcolor=blue
}

\begin{document}

\tableofcontents
\chapter{Introduction aux systèmes automatisés}

% =========================================================
\section{Notion d'automatisme}

\subsection{Introduction}

Dans les installations électriques et mécaniques, certaines opérations doivent être réalisées plusieurs fois, rapidement et avec précision.

Prenons quelques exemples :

\begin{itemize}
    \item allumer automatiquement une lampe lorsqu'il fait sombre ;
    \item démarrer une pompe lorsque le niveau d'eau devient trop bas ;
    \item arrêter un moteur lorsqu'une pièce arrive en fin de course ;
    \item ouvrir automatiquement une porte lorsqu'une personne s'approche ;
    \item déplacer des produits sur un convoyeur ;
    \item remplir automatiquement une bouteille dans une usine.
\end{itemize}

Ces opérations peuvent être réalisées manuellement par une personne. Mais elles peuvent également être exécutées automatiquement grâce à un automatisme.

\subsection{Définition}

Un automatisme est un ensemble de techniques permettant à une machine ou à une installation d'exécuter une ou plusieurs opérations sans intervention humaine permanente, selon des conditions et des règles préalablement définies.

L'automatisme permet donc à une machine de réagir elle-même aux informations qu'elle reçoit.

\paragraph{Exemple :}
Considérons un réservoir d'eau équipé d'une pompe.

Lorsque le niveau d'eau devient trop bas :

\[
\text{Niveau bas} \rightarrow \text{détection} \rightarrow \text{commande} \rightarrow \text{démarrage de la pompe}
\]

Lorsque le réservoir est plein :

\[
\text{Niveau haut} \rightarrow \text{détection} \rightarrow \text{commande} \rightarrow \text{arrêt de la pompe}
\]

L'opérateur n'est donc pas obligé de surveiller continuellement le réservoir.

\subsection{Commande manuelle et commande automatique}

Dans une commande manuelle, l'être humain intervient directement pour provoquer l'action.

\paragraph{Exemple :} pour allumer une lampe, une personne appuie sur un interrupteur.

\[
\text{Personne} \rightarrow \text{interrupteur} \rightarrow \text{lampe}
\]

Dans une commande automatique, une information provenant du système provoque l'action sans intervention permanente de l'opérateur.

\paragraph{Exemple :} un détecteur de luminosité peut commander automatiquement une lampe lorsqu'il fait sombre.

\[
\text{Obscurité} \rightarrow \text{capteur} \rightarrow \text{système de commande} \rightarrow \text{lampe}
\]

La principale différence est donc la suivante :

\begin{table}[h]
\centering
\begin{tabularx}{\textwidth}{|X|X|}
\hline
\textbf{Commande manuelle} & \textbf{Commande automatique} \\
\hline
Nécessite l'intervention de l'opérateur & Fonctionne avec peu ou sans intervention permanente \\
\hline
L'homme prend directement la décision & Le système applique une logique de commande \\
\hline
Exemple : interrupteur d'une lampe & Exemple : éclairage automatique \\
\hline
Adaptée aux opérations simples & Adaptée aux opérations répétitives ou complexes \\
\hline
\end{tabularx}
\caption{Comparaison entre commande manuelle et commande automatique}
\end{table}

\subsection{Automatisme câblé}

Dans un automatisme câblé, le fonctionnement du système est déterminé principalement par la manière dont les composants électriques sont câblés entre eux.

On utilise notamment :

\begin{itemize}
    \item des boutons-poussoirs ;
    \item des relais ;
    \item des contacteurs ;
    \item des temporisateurs ;
    \item des fins de course.
\end{itemize}

\paragraph{Exemple : commande d'un moteur}

\[
\text{Bouton Marche} \rightarrow \text{contacteur} \rightarrow \text{moteur}
\]

Un circuit auxiliaire peut maintenir le contacteur alimenté après le relâchement du bouton Marche. C'est le principe de l'auto-maintien.

Pour modifier profondément le fonctionnement d'un automatisme câblé, il faut souvent modifier son câblage.

\subsection{Automatisme programmé}

Dans un automatisme programmé, une partie importante de la logique de fonctionnement est enregistrée sous forme de programme dans un dispositif électronique.

Dans l'industrie, on utilise notamment l'automate programmable industriel (API), appelé PLC en anglais.

Le principe général est :

\[
\text{Entrées} \rightarrow \text{API} \rightarrow \text{Sorties}
\]

Par exemple :

\[
\text{Bouton Marche} \rightarrow \text{entrée API} \rightarrow \text{programme} \rightarrow \text{sortie API} \rightarrow \text{contacteur} \rightarrow \text{moteur}
\]

Le programme détermine alors dans quelles conditions le moteur doit démarrer ou s'arrêter.

\subsection{Exemple pratique}

Une entreprise possède un convoyeur destiné à transporter des cartons.

Un capteur détecte l'arrivée d'un carton. L'information est envoyée à l'automate. Celui-ci traite l'information selon son programme et commande le moteur du convoyeur.

On obtient :

\[
\text{Capteur} \rightarrow \text{API} \rightarrow \text{contacteur} \rightarrow \text{moteur} \rightarrow \text{convoyeur}
\]

Nous pouvons déjà distinguer trois opérations fondamentales :

\[
\text{Détecter} \rightarrow \text{Traiter} \rightarrow \text{Agir}
\]

Cette logique sera très importante pour comprendre le fonctionnement d'un système automatisé.

\subsection{Importance de l'automatisme}

L'automatisation peut permettre :

\begin{itemize}
    \item d'effectuer des opérations répétitives ;
    \item d'améliorer la régularité et la précision de certaines opérations ;
    \item d'augmenter la cadence de production ;
    \item de surveiller certaines grandeurs automatiquement ;
    \item de réduire certaines interventions manuelles ;
    \item d'améliorer la sécurité lorsqu'elle est correctement conçue ;
    \item de commander plusieurs équipements selon une séquence précise.
\end{itemize}

L'automatisme est aujourd'hui utilisé dans de nombreux domaines : industrie, bâtiments, traitement de l'eau, énergie, transport, mines, agriculture et autres installations techniques.

\subsubsection*{À retenir}

\begin{quote}
Un automatisme permet à une machine ou à une installation d'effectuer des opérations selon des règles déterminées, sans intervention humaine permanente.
\end{quote}

% =========================================================
\section{Définition d'un système automatisé}

\subsection{Définition}

Un système automatisé est un ensemble organisé de composants matériels et, éventuellement, logiciels permettant d'effectuer automatiquement une tâche ou une suite de tâches selon des conditions déterminées.

Il peut recevoir des informations provenant de son environnement, les traiter puis commander différents équipements.

Son fonctionnement général peut être représenté par :

\[
\text{Informations} \rightarrow \text{Traitement} \rightarrow \text{Actions}
\]

Ou plus concrètement :

\[
\text{Capteurs} \rightarrow \text{Système de commande} \rightarrow \text{Actionneurs}
\]

\subsection{Exemple simple : réservoir d'eau}

Considérons un réservoir équipé :

\begin{itemize}
    \item d'un capteur de niveau bas ;
    \item d'un capteur de niveau haut ;
    \item d'un automate programmable ;
    \item d'un contacteur ;
    \item d'une pompe électrique.
\end{itemize}

Lorsque le niveau devient bas, le capteur transmet l'information à l'automate.

L'automate commande alors le contacteur qui alimente la pompe.

Lorsque l'eau atteint le niveau supérieur, le deuxième capteur transmet cette nouvelle information et l'automate arrête la pompe.

Le fonctionnement peut être résumé ainsi :

\[
\text{Niveau bas} \rightarrow \text{Capteur} \rightarrow \text{API} \rightarrow \text{Contacteur} \rightarrow \text{Pompe démarre}
\]

puis :

\[
\text{Niveau haut} \rightarrow \text{Capteur} \rightarrow \text{API} \rightarrow \text{Contacteur} \rightarrow \text{Pompe s'arrête}
\]

\subsection{Les informations d'entrée}

Pour savoir ce qui se passe dans l'installation, le système automatisé doit recevoir des informations d'entrée.

Ces informations peuvent provenir de :

\begin{itemize}
    \item boutons-poussoirs ;
    \item interrupteurs ;
    \item fins de course ;
    \item détecteurs de proximité ;
    \item capteurs de niveau ;
    \item capteurs de température ;
    \item capteurs de pression ;
    \item détecteurs photoélectriques, etc.
\end{itemize}

\paragraph{Exemple :}
Un capteur détecte la présence d'une bouteille sur un convoyeur.

Cette détection constitue une information d'entrée.

\subsection{Traitement des informations}

Les informations reçues doivent ensuite être traitées par le système de commande.

Dans notre cours, nous nous intéresserons principalement à l'automate programmable industriel (API).

L'API possède un programme qui lui indique comment réagir aux différentes situations.

Prenons une règle simple :

\begin{quote}
SI le bouton Marche est activé ET qu'aucun défaut n'est présent, ALORS commander le moteur.
\end{quote}

Le programme de l'automate traduit cette règle en instructions exécutables.

Nous étudierons plus tard comment réaliser cette logique avec le langage Ladder (LD).

\subsection{Les sorties}

Après avoir traité les informations, l'automate commande ses sorties.

Celles-ci peuvent agir directement sur certains petits équipements ou, très souvent, commander des dispositifs intermédiaires qui pilotent les actionneurs.

On peut rencontrer :

\begin{itemize}
    \item voyants ;
    \item relais ;
    \item contacteurs ;
    \item électrovannes ;
    \item avertisseurs sonores.
\end{itemize}

Les actionneurs réalisent l'action physique souhaitée, par exemple :

\begin{itemize}
    \item moteur électrique ;
    \item vérin pneumatique ;
    \item pompe ;
    \item vanne motorisée.
\end{itemize}

\paragraph{Attention :} un moteur industriel de puissance n'est généralement pas alimenté directement par la sortie de l'API. L'API commande typiquement un contacteur ou un variateur, qui assure ensuite la commande électrique du moteur.

\subsection{Exemple d'une porte automatique}

Prenons une porte qui doit s'ouvrir lorsqu'une personne est détectée.

\paragraph{Étape 1 --- Détection}

Le capteur détecte une personne.

\paragraph{Étape 2 --- Information}

Le capteur transmet un signal au système de commande.

\paragraph{Étape 3 --- Traitement}

Le système vérifie les conditions prévues par sa logique.

\paragraph{Étape 4 --- Commande}

Il commande l'équipement chargé d'actionner la porte.

\paragraph{Étape 5 --- Action}

Le moteur entraîne le mécanisme et la porte s'ouvre.

Nous retrouvons donc :

\[
\text{CAPTEUR} \rightarrow \text{TRAITEMENT} \rightarrow \text{COMMANDE} \rightarrow \text{ACTIONNEUR} \rightarrow \text{ACTION}
\]

\subsection{Système automatisé et opérateur humain}

Un système automatisé ne signifie pas nécessairement que l'être humain disparaît complètement.

L'opérateur peut toujours intervenir pour :

\begin{itemize}
    \item mettre l'installation en marche ;
    \item sélectionner un mode de fonctionnement ;
    \item arrêter la machine ;
    \item surveiller son fonctionnement ;
    \item acquitter certaines alarmes ;
    \item effectuer la maintenance ;
    \item intervenir en cas de panne.
\end{itemize}

L'automatisation consiste surtout à confier au système les opérations qui ont été prévues lors de sa conception.

\subsection{Exemple industriel récapitulatif}

Considérons une machine qui remplit automatiquement des bouteilles.

La bouteille arrive devant le poste de remplissage.

\[
\text{Capteur de présence} \rightarrow \text{API} \rightarrow \text{électrovanne} \rightarrow \text{remplissage}
\]

Lorsque les conditions prévues sont satisfaites, l'automate commande l'ouverture de l'électrovanne. Après la quantité ou la durée prévue par la conception du système, l'électrovanne est refermée et la séquence peut continuer.

Ce simple exemple contient déjà les principaux éléments que nous étudierons :

\[
\text{capteurs} + \text{automate} + \text{programme} + \text{sorties} + \text{actionneurs} + \text{machine}.
\]

\subsubsection*{À retenir}

\begin{quote}
Un système automatisé reçoit des informations, les traite selon une logique déterminée et commande des équipements afin de réaliser automatiquement une tâche.
\end{quote}

\paragraph{Schéma fondamental :}

\[
\text{ENTRÉES} \rightarrow \text{TRAITEMENT} \rightarrow \text{SORTIES}
\]

\[
\text{Capteurs} \rightarrow \text{API} \rightarrow \text{Actionneurs}
\]

% =========================================================
\section{Objectifs de l'automatisation}

\subsection{Introduction}

Automatiser une machine ou une installation n'est pas une fin en soi. Une entreprise automatise pour répondre à des besoins précis : améliorer la production, garantir la qualité, réduire les coûts, protéger les personnes ou surveiller un procédé.

L'automatisation doit donc être justifiée par un ou plusieurs objectifs clairement définis dans le cahier des charges.

\subsection{Pourquoi automatiser ?}

On automatise principalement pour :

\begin{itemize}
    \item effectuer des opérations répétitives sans fatigue ;
    \item obtenir une production régulière et conforme ;
    \item augmenter la cadence de production ;
    \item réduire les erreurs humaines ;
    \item surveiller en permanence certaines grandeurs ;
    \item protéger les opérateurs contre les tâches dangereuses ;
    \item commander plusieurs équipements selon une séquence précise ;
    \item réduire certaines consommations d'énergie ou de matières.
\end{itemize}

\subsection{Objectifs techniques}

Sur le plan technique, l'automatisation permet :

\begin{itemize}
    \item la répétitivité des opérations ;
    \item la précision des actions ;
    \item la rapidité d'exécution ;
    \item la coordination de plusieurs actionneurs ;
    \item la surveillance continue du système ;
    \item la gestion de séquences complexes.
\end{itemize}

\paragraph{Exemple :} une ligne de remplissage peut remplir, peser, boucher et étiqueter des bouteilles dans un ordre précis, sans intervention permanente de l'opérateur.

\subsection{Objectifs économiques}

L'automatisation cherche souvent à améliorer la rentabilité de l'installation.

Elle peut permettre de :

\begin{itemize}
    \item réduire le temps de cycle ;
    \item diminuer les rebuts et les pertes de matière ;
    \item réduire certaines interventions manuelles ;
    \item augmenter le rendement des machines ;
    \item améliorer la disponibilité de l'installation ;
    \item produire avec une qualité plus constante.
\end{itemize}

Cependant, l'automatisation représente aussi un investissement : achat du matériel, câblage, programmation, formation du personnel et maintenance.

\subsection{Objectifs liés à la sécurité et aux conditions de travail}

Un système automatisé bien conçu peut :

\begin{itemize}
    \item éloigner l'opérateur des zones dangereuses ;
    \item supprimer certaines tâches pénibles ou répétitives ;
    \item détecter plus rapidement un défaut ou une anomalie ;
    \item déclencher un arrêt d'urgence ou une alarme ;
    \item imposer une séquence de fonctionnement plus sûre.
\end{itemize}

L'opérateur devient alors davantage un superviseur qu'un exécutant direct.

\paragraph{Attention :} automatiser ne signifie pas supprimer toute surveillance humaine. La sécurité doit être conçue dès le départ et entretenue pendant toute la durée de vie de l'installation.

\subsection{Objectifs de qualité et de traçabilité}

L'automatisation permet d'obtenir :

\begin{itemize}
    \item une qualité régulière ;
    \item le respect strict d'un cahier des charges ;
    \item un contrôle automatique de certaines caractéristiques ;
    \item l'enregistrement de données de production ;
    \item la génération d'alarmes en cas de dérive.
\end{itemize}

Dans certaines industries, il est important de savoir quand, comment et avec quels paramètres un produit a été fabriqué. C'est le rôle de la traçabilité.

\subsection{Objectifs de flexibilité}

Un automatisme programmé, comme celui réalisé avec un API, peut être modifié plus facilement qu'un automatisme entièrement câblé.

Il devient alors possible de :

\begin{itemize}
    \item changer un format de production ;
    \item modifier une recette ou une séquence ;
    \item adapter la cadence ;
    \item ajouter une nouvelle étape ;
    \item passer d'un mode manuel à un mode automatique.
\end{itemize}

Cette flexibilité est très importante dans les industries où les produits ou les demandes changent souvent.

\subsection{Exemples industriels}

\begin{itemize}
    \item \textbf{Traitement de l'eau :} démarrage et arrêt automatiques des pompes selon le niveau des réservoirs.
    \item \textbf{Convoyeur :} détection des colis et commande du moteur.
    \item \textbf{Remplissage de cuve :} ouverture d'une électrovanne jusqu'à un niveau ou un volume déterminé.
    \item \textbf{Porte automatique :} détection d'une personne et commande du moteur d'ouverture.
    \item \textbf{Tri de pièces :} identification et orientation des pièces selon leur type.
    \item \textbf{Ascenseur simplifié :} gestion des appels, des niveaux et des sécurités.
\end{itemize}

\subsection{Avantages et limites de l'automatisation}

\begin{table}[h]
\centering
\begin{tabularx}{\textwidth}{|X|X|}
\hline
\textbf{Avantages} & \textbf{Limites / points de vigilance} \\
\hline
Augmentation de la productivité & Coût d'investissement initial \\
\hline
Qualité plus régulière & Complexité de conception et de maintenance \\
\hline
Réduction de certaines tâches pénibles & Besoin de personnel qualifié \\
\hline
Meilleure sécurité lorsqu'elle est bien conçue & Dépendance à l'énergie et à l'électronique \\
\hline
Surveillance continue du procédé & Risque d'obsolescence du matériel \\
\hline
Flexibilité grâce à la programmation & Nécessité de formation et de documentation \\
\hline
\end{tabularx}
\caption{Avantages et limites de l'automatisation}
\end{table}

\subsection{Conclusion}

Automatiser consiste à confier à un système technique des opérations qui étaient auparavant réalisées manuellement ou avec une aide limitée.

Les objectifs peuvent être techniques, économiques, humains ou liés à la qualité. Ils doivent être clairement définis avant de choisir les capteurs, les actionneurs, l'automate et le programme.

\subsubsection*{À retenir}

\begin{quote}
L'automatisation vise à améliorer la production, la qualité, la sécurité et la régularité d'un procédé, tout en tenant compte des coûts, de la maintenance et des compétences nécessaires.
\end{quote}

\paragraph{Schéma fondamental :}

\[
\text{Besoin} \rightarrow \text{Objectifs} \rightarrow \text{Choix techniques} \rightarrow \text{Réalisation}
\]
% =========================================================
\section{Structure générale d'un système automatisé}

\subsection{Introduction}

Pour comprendre le fonctionnement d'un système automatisé, il ne suffit pas de connaître ses composants. Il faut aussi savoir comment ils sont organisés et comment ils échangent des informations.

Un système automatisé est généralement structuré en deux grandes parties :

\begin{itemize}
    \item la \textbf{partie opérative} (PO), qui agit sur la matière d'œuvre ;
    \item la \textbf{partie commande} (PC), qui élabore les ordres.
\end{itemize}

À cela s'ajoutent les interfaces de dialogue et les chaînes d'information et d'énergie qui relient les différents éléments.

\subsection{Partie opérative et partie commande}

\subsubsection{La partie opérative (PO)}

La partie opérative est l'ensemble des éléments qui réalisent physiquement le travail sur la matière d'œuvre.

Elle comprend notamment :

\begin{itemize}
    \item les \textbf{actionneurs} : moteurs, vérins, électrovannes, etc. ;
    \item les \textbf{préactionneurs} : contacteurs, relais, distributeurs, variateurs ;
    \item les \textbf{effecteurs} : outils, pinces, convoyeurs, etc. ;
    \item les \textbf{capteurs} qui informent sur l'état de la PO (position, présence, niveau, etc.).
\end{itemize}

\paragraph{Exemple :} dans un convoyeur, la partie opérative comprend le moteur, le réducteur, le tapis et les capteurs de présence.

\subsubsection{La partie commande (PC)}

La partie commande est l'ensemble des éléments qui traitent les informations et élaborent les ordres destinés à la partie opérative.

Elle peut être réalisée par :

\begin{itemize}
    \item des relais et contacteurs (automatisme câblé) ;
    \item un automate programmable industriel (API) ;
    \item un calculateur industriel ou un microcontrôleur.
\end{itemize}

La partie commande reçoit les informations des capteurs, applique la logique du programme et commande les sorties.

\subsubsection{Relation entre PO et PC}

La partie commande et la partie opérative échangent en permanence :

\begin{itemize}
    \item la PC envoie des \textbf{ordres} à la PO ;
    \item la PO renvoie des \textbf{comptes rendus} à la PC.
\end{itemize}

On peut représenter cet échange par :

\[
\text{PC} \xrightarrow{\text{ordres}} \text{PO}
\qquad \text{et} \qquad
\text{PO} \xrightarrow{\text{comptes rendus}} \text{PC}
\]

\subsection{Chaîne d'information et chaîne d'énergie}

Dans un système automatisé, on distingue souvent deux chaînes fonctionnelles.

\subsubsection{La chaîne d'information}

Elle concerne l'acquisition, le traitement et la communication des informations.

Ses éléments principaux sont :

\begin{itemize}
    \item \textbf{Acquérir} : capteurs, boutons, détecteurs ;
    \item \textbf{Traiter} : API, relais, calculateur ;
    \item \textbf{Communiquer} : interfaces, réseaux industriels, IHM.
\end{itemize}

\paragraph{Exemple :} un capteur de niveau envoie une information à l'API. L'API la traite et affiche l'état sur une IHM.

\subsubsection{La chaîne d'énergie}

Elle concerne l'alimentation, la distribution et la conversion de l'énergie nécessaire aux actionneurs.

Ses éléments principaux sont :

\begin{itemize}
    \item \textbf{Alimenter} : réseau électrique, air comprimé, hydraulique ;
    \item \textbf{Distribuer} : contacteurs, disjoncteurs, distributeurs ;
    \item \textbf{Convertir} : moteurs, vérins, résistances ;
    \item \textbf{Transmettre} : arbres, courroies, engrenages, tuyaux.
\end{itemize}

\paragraph{Exemple :} le réseau électrique alimente un contacteur qui alimente un moteur. Le moteur convertit l'énergie électrique en énergie mécanique.

\subsection{Schéma général d'un système automatisé}

On peut représenter la structure générale de la manière suivante :

\[
\begin{array}{c}
\text{Opérateur} \\
\downarrow \\
\text{IHM / Pupitre} \\
\downarrow \\
\text{Partie Commande (API)} \\
\swarrow \qquad \searrow \\
\text{Ordres} \qquad \text{Comptes rendus} \\
\downarrow \qquad \uparrow \\
\text{Partie Opérative} \\
\downarrow \\
\text{Action sur la matière d'œuvre}
\end{array}
\]

Une autre représentation simplifiée est :

\[
\text{Capteurs} \rightarrow \text{Partie Commande} \rightarrow \text{Actionneurs}
\]

\[
\text{Chaîne d'information} \qquad \text{Chaîne d'énergie}
\]

\subsection{Exemple : réservoir d'eau}

Reprenons l'exemple du réservoir d'eau.

\begin{itemize}
    \item \textbf{Partie opérative} : pompe, contacteur, réservoir, capteurs de niveau.
    \item \textbf{Partie commande} : API.
    \item \textbf{Chaîne d'information} : capteurs de niveau $\rightarrow$ API.
    \item \textbf{Chaîne d'énergie} : API $\rightarrow$ contacteur $\rightarrow$ pompe.
\end{itemize}

Le cycle est le suivant :

\[
\text{Niveau bas} \rightarrow \text{Capteur} \rightarrow \text{API} \rightarrow \text{Contacteur} \rightarrow \text{Pompe}
\]

\[
\text{Niveau haut} \rightarrow \text{Capteur} \rightarrow \text{API} \rightarrow \text{Contacteur} \rightarrow \text{Arrêt pompe}
\]

\subsection{Exemple : convoyeur}

\begin{itemize}
    \item \textbf{Partie opérative} : moteur, tapis, capteur de présence.
    \item \textbf{Partie commande} : API.
    \item \textbf{Chaîne d'information} : capteur $\rightarrow$ API.
    \item \textbf{Chaîne d'énergie} : API $\rightarrow$ contacteur $\rightarrow$ moteur.
\end{itemize}

\[
\text{Carton détecté} \rightarrow \text{API} \rightarrow \text{Contacteur} \rightarrow \text{Moteur} \rightarrow \text{Convoyeur avance}
\]

\subsection{Place de l'opérateur}

L'opérateur n'est pas exclu du système. Il intervient par l'intermédiaire d'une interface homme-machine (IHM) ou d'un pupitre.

Il peut :

\begin{itemize}
    \item mettre en marche ou arrêter l'installation ;
    \item choisir un mode de fonctionnement ;
    \item acquitter les alarmes ;
    \item surveiller les paramètres ;
    \item intervenir en maintenance.
\end{itemize}

\subsection{Tableau récapitulatif}

\begin{table}[h]
\centering
\begin{tabularx}{\textwidth}{|X|X|}
\hline
\textbf{Partie commande (PC)} & \textbf{Partie opérative (PO)} \\
\hline
Traite les informations & Agit sur la matière d'œuvre \\
\hline
API, relais, calculateur & Moteurs, vérins, contacteurs \\
\hline
Reçoit les comptes rendus & Reçoit les ordres \\
\hline
Élabore les ordres & Exécute les actions \\
\hline
\end{tabularx}
\caption{Répartition des rôles entre PC et PO}
\end{table}

\subsection{Conclusion}

La structure d'un système automatisé repose sur une séparation claire entre la partie commande et la partie opérative. Les informations circulent dans la chaîne d'information, tandis que l'énergie est distribuée dans la chaîne d'énergie.

Cette organisation permet de comprendre comment un API peut commander une machine, tout en restant en relation avec les capteurs, les actionneurs et l'opérateur.

\subsubsection*{À retenir}

\begin{quote}
Un système automatisé comprend une partie commande qui traite les informations et une partie opérative qui exécute les actions. Les deux communiquent par des ordres et des comptes rendus.
\end{quote}

\paragraph{Schéma fondamental :}

\[
\text{Capteurs} \rightarrow \text{Partie Commande} \rightarrow \text{Actionneurs}
\]

\[
\text{Chaîne d'information} \qquad \text{Chaîne d'énergie}
\]
% =========================================================
\section{Partie commande et partie opérative}

\subsection{Introduction}

Dans un système automatisé, la \textbf{partie opérative} (PO) exécute les actions physiques, tandis que la \textbf{partie commande} (PC) décide de ces actions à partir des informations reçues.

Cette séparation est fondamentale : elle permet de distinguer clairement ce qui \textbf{décide} de ce qui \textbf{agit}.

\subsection{La partie opérative (PO)}

\subsubsection{Définition}

La partie opérative est l'ensemble des composants qui réalisent physiquement le travail sur la matière d'œuvre, sous l'ordre de la partie commande.

Elle agit sur la matière, la déplace, la transforme, la mesure ou la contrôle.

\subsubsection{Composants de la partie opérative}

On distingue généralement :

\begin{itemize}
    \item les \textbf{actionneurs} : moteurs, vérins, électroaimants ;
    \item les \textbf{préactionneurs} : contacteurs, relais, distributeurs, variateurs ;
    \item les \textbf{effecteurs} : outils, pinces, buses, tapis ;
    \item les \textbf{capteurs} : fins de course, détecteurs de proximité, capteurs de niveau ;
    \item les \textbf{organes de transmission} : arbres, courroies, engrenages, tuyaux.
\end{itemize}

\subsubsection{Rôle de la partie opérative}

La PO a pour fonction de :

\begin{itemize}
    \item recevoir les ordres de la PC ;
    \item convertir une énergie (électrique, pneumatique, hydraulique) en action mécanique ;
    \item agir sur la matière d'œuvre ;
    \item informer la PC de son état par l'intermédiaire des capteurs.
\end{itemize}

\paragraph{Exemple :} dans une porte automatique, la PO comprend le moteur, le réducteur, le bras mécanique et les capteurs de présence.

\subsection{La partie commande (PC)}

\subsubsection{Définition}

La partie commande est l'ensemble des composants qui traitent les informations et élaborent les ordres destinés à la partie opérative.

Elle constitue le « cerveau » du système automatisé.

\subsubsection{Composants de la partie commande}

Selon le type d'automatisme, la PC peut être réalisée par :

\begin{itemize}
    \item des \textbf{relais} et \textbf{contacteurs} (automatisme câblé) ;
    \item un \textbf{automate programmable industriel (API)} ;
    \item un \textbf{microcontrôleur} ou un calculateur industriel ;
    \item des \textbf{modules logiques} programmables.
\end{itemize}

\subsubsection{Rôle de la partie commande}

La PC a pour fonction de :

\begin{itemize}
    \item acquérir les informations d'entrée ;
    \item traiter ces informations selon une logique définie ;
    \item élaborer les ordres de sortie ;
    \item transmettre ces ordres à la PO ;
    \item gérer les sécurités, les alarmes et les modes de fonctionnement.
\end{itemize}

\paragraph{Exemple :} dans une porte automatique, la PC est un API qui reçoit l'information du capteur de présence et commande le moteur d'ouverture.

\subsection{Dialogue entre PC et PO}

La PC et la PO communiquent en permanence par deux flux d'informations :

\begin{itemize}
    \item la PC envoie des \textbf{ordres} à la PO ;
    \item la PO renvoie des \textbf{comptes rendus} à la PC.
\end{itemize}

On peut représenter ce dialogue par :

\[
\text{PC} \xrightarrow{\text{ordres}} \text{PO}
\qquad \text{et} \qquad
\text{PO} \xrightarrow{\text{comptes rendus}} \text{PC}
\]

\paragraph{Exemple :}

\begin{itemize}
    \item Ordre : « démarrer le moteur du convoyeur ».
    \item Compte rendu : « le convoyeur est en marche » ou « le carton est détecté ».
\end{itemize}

\subsection{Les informations échangées}

\subsubsection{Les ordres}

Les ordres sont élaborés par la PC. Ils peuvent être :

\begin{itemize}
    \item \textbf{TOR} (Tout Ou Rien) : marcher/arrêter, ouvrir/fermer ;
    \item \textbf{analogiques} : consigne de vitesse, de position, de température ;
    \item \textbf{numériques} : envoi d'un code, d'un compteur, d'une recette.
\end{itemize}

\subsubsection{Les comptes rendus}

Les comptes rendus sont fournis par les capteurs de la PO. Ils peuvent être :

\begin{itemize}
    \item \textbf{TOR} : présence, absence, fin de course ;
    \item \textbf{analogiques} : niveau, pression, température, débit ;
    \item \textbf{numériques} : position d'un codeur, comptage.
\end{itemize}

\subsection{Exemple complet : remplissage de cuve}

Considérons une cuve à remplir automatiquement.

\paragraph{Partie opérative :}
\begin{itemize}
    \item électrovanne de remplissage ;
    \item pompe d'alimentation ;
    \item capteurs de niveau bas et haut ;
    \item cuve.
\end{itemize}

\paragraph{Partie commande :}
\begin{itemize}
    \item API.
\end{itemize}

\paragraph{Dialogue :}

\[
\text{Niveau bas détecté} \rightarrow \text{API} \rightarrow \text{Ordre d'ouverture de l'électrovanne}
\]

\[
\text{Niveau haut détecté} \rightarrow \text{API} \rightarrow \text{Ordre de fermeture de l'électrovanne}
\]

\paragraph{Analyse :}
\begin{itemize}
    \item Le capteur de niveau bas est un \textbf{compte rendu} envoyé par la PO.
    \item L'API élabore un \textbf{ordre} d'ouverture ou de fermeture.
    \item L'électrovanne est un \textbf{actionneur} de la PO.
\end{itemize}

\subsection{Exemple complet : convoyeur}

\paragraph{Partie opérative :}
\begin{itemize}
    \item moteur du convoyeur ;
    \item tapis ;
    \item capteur de présence de carton ;
    \item contacteur.
\end{itemize}

\paragraph{Partie commande :}
\begin{itemize}
    \item API.
\end{itemize}

\paragraph{Dialogue :}

\[
\text{Carton détecté} \rightarrow \text{API} \rightarrow \text{Contacteur} \rightarrow \text{Moteur}
\]

\[
\text{Fin de convoyage} \rightarrow \text{API} \rightarrow \text{Arrêt du moteur}
\]

\subsection{Place de l'IHM}

L'IHM (Interface Homme-Machine) permet à l'opérateur de dialoguer avec la partie commande.

Elle peut :

\begin{itemize}
    \item envoyer des ordres à la PC (marche, arrêt, choix de mode) ;
    \item afficher les comptes rendus (états, alarmes, mesures) ;
    \item visualiser le fonctionnement du système.
\end{itemize}

Elle est reliée à la PC et non directement à la PO.

\subsection{Tableau récapitulatif}

\begin{table}[h]
\centering
\begin{tabularx}{\textwidth}{|X|X|}
\hline
\textbf{Partie commande (PC)} & \textbf{Partie opérative (PO)} \\
\hline
Traite les informations & Exécute les actions \\
\hline
API, relais, calculateur & Moteurs, vérins, contacteurs \\
\hline
Reçoit les comptes rendus & Fournit les comptes rendus \\
\hline
Élabore les ordres & Reçoit et exécute les ordres \\
\hline
Décide & Agit \\
\hline
\end{tabularx}
\caption{Répartition des rôles entre PC et PO}
\end{table}

\subsection{Conclusion}

La distinction entre partie commande et partie opérative est essentielle pour comprendre et concevoir un système automatisé.

La PC \textbf{décide}, la PO \textbf{agit}. Elles communiquent par des ordres et des comptes rendus. L'IHM permet à l'opérateur d'intervenir sans agir directement sur la PO.

\subsubsection*{À retenir}

\begin{quote}
La partie commande élabore les ordres à partir des informations reçues. La partie opérative exécute ces ordres sur la matière d'œuvre et renvoie des comptes rendus. L'IHM permet à l'opérateur de dialoguer avec la partie commande.
\end{quote}

\paragraph{Schéma fondamental :}

\[
\text{IHM} \rightarrow \text{Partie Commande} \rightarrow \text{Partie Opérative} \rightarrow \text{Matière d'œuvre}
\]

\[
\text{Ordres} \ \downarrow \qquad \uparrow \ \text{Comptes rendus}
\]
% =========================================================
\section{Chaîne d'information et chaîne d'énergie}

\subsection{Introduction}

Pour fonctionner, un système automatisé a besoin de deux choses essentielles :

\begin{itemize}
    \item des \textbf{informations}, pour savoir ce qui se passe et décider quoi faire ;
    \item de l'\textbf{énergie}, pour agir physiquement sur la matière d'œuvre.
\end{itemize}

Ces deux besoins correspondent à deux chaînes fonctionnelles complémentaires :

\begin{itemize}
    \item la \textbf{chaîne d'information} ;
    \item la \textbf{chaîne d'énergie}.
\end{itemize}

Comprendre ces deux chaînes permet de savoir où circule l'information et où circule l'énergie dans une installation automatisée.

\subsection{La chaîne d'information}

\subsubsection{Définition}

La chaîne d'information regroupe l'ensemble des éléments qui permettent d'\textbf{acquérir}, de \textbf{traiter} et de \textbf{communiquer} des informations dans un système automatisé.

Elle ne transporte pas d'énergie de puissance : elle transporte des \textbf{signaux}.

\subsubsection{Fonctions de la chaîne d'information}

La chaîne d'information assure trois fonctions principales :

\begin{itemize}
    \item \textbf{Acquérir} : recueillir les informations provenant des capteurs, boutons, détecteurs ;
    \item \textbf{Traiter} : analyser ces informations selon la logique du programme ;
    \item \textbf{Communiquer} : transmettre les informations vers les interfaces, l'IHM ou d'autres automates.
\end{itemize}

\subsubsection{Éléments de la chaîne d'information}

On retrouve notamment :

\begin{itemize}
    \item les \textbf{capteurs} (fins de course, détecteurs de proximité, capteurs de niveau, etc.) ;
    \item les \textbf{boutons-poussoirs} et \textbf{interrupteurs} ;
    \item l'\textbf{automate programmable industriel (API)} ;
    \item les \textbf{interfaces de communication} (réseaux, bus de terrain) ;
    \item l'\textbf{IHM} (Interface Homme-Machine).
\end{itemize}

\paragraph{Exemple :} un capteur de niveau envoie une information à l'API. L'API la traite et affiche l'état sur une IHM.

\subsection{La chaîne d'énergie}

\subsubsection{Définition}

La chaîne d'énergie regroupe l'ensemble des éléments qui permettent d'\textbf{alimenter}, de \textbf{distribuer}, de \textbf{convertir} et de \textbf{transmettre} l'énergie nécessaire aux actionneurs.

Elle transporte de l'\textbf{énergie de puissance}.

\subsubsection{Fonctions de la chaîne d'énergie}

La chaîne d'énergie assure quatre fonctions principales :

\begin{itemize}
    \item \textbf{Alimenter} : fournir l'énergie au système (réseau électrique, air comprimé, hydraulique) ;
    \item \textbf{Distribuer} : diriger l'énergie vers les actionneurs (contacteurs, disjoncteurs, distributeurs) ;
    \item \textbf{Convertir} : transformer l'énergie en travail mécanique (moteurs, vérins) ;
    \item \textbf{Transmettre} : transmettre le mouvement à l'effecteur (arbres, courroies, engrenages).
\end{itemize}

\subsubsection{Éléments de la chaîne d'énergie}

On retrouve notamment :

\begin{itemize}
    \item l'\textbf{alimentation} électrique ou pneumatique ;
    \item les \textbf{contacteurs} et \textbf{disjoncteurs} ;
    \item les \textbf{moteurs} électriques ;
    \item les \textbf{vérins} pneumatiques ou hydrauliques ;
    \item les \textbf{transmissions mécaniques}.
\end{itemize}

\paragraph{Exemple :} le réseau électrique alimente un contacteur qui alimente un moteur. Le moteur convertit l'énergie électrique en énergie mécanique pour entraîner le convoyeur.

\subsection{Comparaison entre les deux chaînes}

\begin{table}[h]
\centering
\begin{tabularx}{\textwidth}{|X|X|}
\hline
\textbf{Chaîne d'information} & \textbf{Chaîne d'énergie} \\
\hline
Transporte des \textbf{signaux} & Transporte de l'\textbf{énergie} \\
\hline
Faible puissance & Puissance importante \\
\hline
Capteurs, API, IHM & Alimentation, contacteurs, moteurs \\
\hline
Acquérir, traiter, communiquer & Alimenter, distribuer, convertir, transmettre \\
\hline
\end{tabularx}
\caption{Comparaison entre chaîne d'information et chaîne d'énergie}
\end{table}

\subsection{Relation entre les deux chaînes}

Les deux chaînes ne fonctionnent pas séparément : elles sont \textbf{reliées} par la partie commande.

La chaîne d'information :

\begin{itemize}
    \item reçoit les informations des capteurs ;
    \item les traite dans l'API ;
    \item élabore les ordres.
\end{itemize}

La chaîne d'énergie :

\begin{itemize}
    \item reçoit les ordres de la partie commande ;
    \item distribue l'énergie vers les actionneurs ;
    \item réalise l'action sur la matière d'œuvre.
\end{itemize}

On peut représenter cette relation par :

\[
\text{Capteurs} \rightarrow \text{Chaîne d'information} \rightarrow \text{Ordres}
\]

\[
\text{Ordres} \rightarrow \text{Chaîne d'énergie} \rightarrow \text{Actionneurs} \rightarrow \text{Action}
\]

\subsection{Représentation graphique}

Une représentation classique d'un système automatisé sous forme de deux chaînes est la suivante :

\[
\begin{array}{ccccccc}
\text{Acquérir} & \rightarrow & \text{Traiter} & \rightarrow & \text{Communiquer} & & \\
\text{(Capteurs)} & & \text{(API)} & & \text{(IHM, réseau)} & & \\
& & \downarrow & & & & \\
& & \text{Ordres} & & & & \\
& & \downarrow & & & & \\
\text{Alimenter} & \rightarrow & \text{Distribuer} & \rightarrow & \text{Convertir} & \rightarrow & \text{Transmettre} \\
\text{(Réseau)} & & \text{(Contacteurs)} & & \text{(Moteurs)} & & \text{(Mécanique)}
\end{array}
\]

\subsection{Exemple : réservoir d'eau}

Reprenons l'exemple du réservoir d'eau.

\paragraph{Chaîne d'information :}
\begin{itemize}
    \item capteur de niveau bas ;
    \item capteur de niveau haut ;
    \item API.
\end{itemize}

\paragraph{Chaîne d'énergie :}
\begin{itemize}
    \item réseau électrique ;
    \item contacteur ;
    \item pompe électrique ;
    \item tuyauterie.
\end{itemize}

\paragraph{Fonctionnement :}

\[
\text{Niveau bas} \rightarrow \text{Capteur} \rightarrow \text{API} \rightarrow \text{Contacteur} \rightarrow \text{Pompe}
\]

\[
\text{Niveau haut} \rightarrow \text{Capteur} \rightarrow \text{API} \rightarrow \text{Contacteur} \rightarrow \text{Arrêt pompe}
\]

\subsection{Exemple : convoyeur}

\paragraph{Chaîne d'information :}
\begin{itemize}
    \item capteur de présence de carton ;
    \item API ;
    \item IHM.
\end{itemize}

\paragraph{Chaîne d'énergie :}
\begin{itemize}
    \item réseau électrique ;
    \item contacteur ;
    \item moteur ;
    \item tapis du convoyeur.
\end{itemize}

\paragraph{Fonctionnement :}

\[
\text{Carton détecté} \rightarrow \text{API} \rightarrow \text{Contacteur} \rightarrow \text{Moteur} \rightarrow \text{Convoyeur avance}
\]

\subsection{Exemple : porte automatique}

\paragraph{Chaîne d'information :}
\begin{itemize}
    \item capteur de présence ;
    \item API.
\end{itemize}

\paragraph{Chaîne d'énergie :}
\begin{itemize}
    \item réseau électrique ;
    \item variateur ou contacteur ;
    \item moteur ;
    \item mécanisme de la porte.
\end{itemize}

\paragraph{Fonctionnement :}

\[
\text{Personne détectée} \rightarrow \text{API} \rightarrow \text{Moteur} \rightarrow \text{Porte s'ouvre}
\]

\subsection{Importance de cette distinction}

Distinguer la chaîne d'information de la chaîne d'énergie permet :

\begin{itemize}
    \item de comprendre le rôle de chaque composant ;
    \item de localiser une panne plus facilement ;
    \item de choisir correctement les capteurs et les actionneurs ;
    \item de dimensionner les câbles et les protections ;
    \item de respecter les règles de sécurité.
\end{itemize}

\paragraph{Attention :} dans un schéma, on ne mélange pas les deux chaînes. Les capteurs et l'API appartiennent à la chaîne d'information ; les contacteurs, moteurs et vérins appartiennent à la chaîne d'énergie.

\subsection{Conclusion}

Tout système automatisé comporte deux chaînes complémentaires : la chaîne d'information, qui transporte les signaux, et la chaîne d'énergie, qui transporte la puissance.

La partie commande fait le lien entre les deux : elle reçoit les informations de la chaîne d'information et envoie les ordres à la chaîne d'énergie.

\subsubsection*{À retenir}

\begin{quote}
La chaîne d'information acquiert, traite et communique les signaux. La chaîne d'énergie alimente, distribue, convertit et transmet l'énergie. La partie commande relie ces deux chaînes.
\end{quote}

\paragraph{Schéma fondamental :}

\[
\text{Capteurs} \rightarrow \text{Chaîne d'information} \rightarrow \text{Ordres}
\]

\[
\text{Ordres} \rightarrow \text{Chaîne d'énergie} \rightarrow \text{Actionneurs} \rightarrow \text{Matière d'œuvre}
\]
% =========================================================
\chapter{Éléments d'un système automatisé}

\section{Organes d'entrée}

\subsection{Introduction}

Pour fonctionner correctement, un système automatisé doit d'abord \textbf{savoir ce qui se passe} dans son environnement et dans sa partie opérative.

Les \textbf{organes d'entrée} sont les composants qui fournissent ces informations à la partie commande. Ils constituent le premier maillon de la chaîne d'information.

Sans organes d'entrée, l'API ne peut ni détecter une situation, ni décider d'une action.

\subsection{Définition}

Un organe d'entrée est un composant qui transforme une information physique (position, présence, niveau, température, pression, etc.) en un signal exploitable par la partie commande.

Ce signal peut être :

\begin{itemize}
    \item \textbf{TOR} (Tout Ou Rien) : présence/absence, ouvert/fermé ;
    \item \textbf{analogique} : tension ou courant proportionnel à une grandeur ;
    \item \textbf{numérique} : code transmis sous forme de mots binaires.
\end{itemize}

\subsection{Rôle des organes d'entrée}

Les organes d'entrée ont pour fonction de :

\begin{itemize}
    \item détecter un état ou un événement ;
    \item transmettre cette information à la partie commande ;
    \item permettre à l'API de décider de l'action à exécuter ;
    \item renseigner l'opérateur sur l'état du système.
\end{itemize}

Ils constituent les \textbf{yeux} et les \textbf{oreilles} du système automatisé.

\subsection{Classification des organes d'entrée}

On peut classer les organes d'entrée en plusieurs catégories.

\subsubsection{Organes de commande manuelle}

Ils sont actionnés par l'opérateur.

\begin{itemize}
    \item Boutons-poussoirs (Marche, Arrêt, Arrêt d'urgence) ;
    \item Interrupteurs et commutateurs ;
    \item Sélecteurs (mode manuel / automatique) ;
    \item Boutons rotatifs ou à clé.
\end{itemize}

\subsubsection{Capteurs de position et de présence}

Ils détectent la présence ou la position d'un objet ou d'un organe.

\begin{itemize}
    \item Fins de course mécaniques ;
    \item Détecteurs de proximité inductifs ;
    \item Détecteurs de proximité capacitifs ;
    \item Détecteurs photoélectriques ;
    \item Détecteurs à ultrasons.
\end{itemize}

\subsubsection{Capteurs de grandeur physique}

Ils mesurent une grandeur continue.

\begin{itemize}
    \item Capteurs de niveau ;
    \item Capteurs de température ;
    \item Capteurs de pression ;
    \item Capteurs de débit ;
    \item Capteurs de position analogiques ;
    \item Codeurs.
\end{itemize}

\subsubsection{Organes de sécurité}

Ils assurent la protection des personnes et des machines.

\begin{itemize}
    \item Boutons d'arrêt d'urgence ;
    \item Barrières immatérielles ;
    \item Capteurs de porte ;
    \item Relais de sécurité.
\end{itemize}

\subsection{Les boutons-poussoirs}

Un bouton-poussoir est un organe de commande manuelle qui ferme ou ouvre un contact lorsqu'on appuie dessus.

On distingue :

\begin{itemize}
    \item \textbf{bouton à contact NO} (Normalement Ouvert) : le contact se ferme quand on appuie ;
    \item \textbf{bouton à contact NC} (Normalement Fermé) : le contact s'ouvre quand on appuie.
\end{itemize}

\paragraph{Exemple :}

\begin{itemize}
    \item Bouton \textbf{Marche} : contact NO ;
    \item Bouton \textbf{Arrêt} : contact NC ;
    \item Bouton \textbf{Arrêt d'urgence} : contact NC à verrouillage.
\end{itemize}

\subsection{Les interrupteurs et commutateurs}

Un interrupteur ou commutateur possède deux positions stables. Il conserve son état après l'action de l'opérateur.

\paragraph{Exemple :} un sélecteur « Manuel / Automatique » qui reste dans la position choisie.

\subsection{Les capteurs et détecteurs}

\subsubsection{Fin de course}

Un fin de course est un capteur mécanique qui détecte la position extrême d'un organe mobile.

\paragraph{Exemple :} détection de la position d'un vérin complètement sorti.

\subsubsection{Détecteur de proximité inductif}

Il détecte la présence d'un objet métallique sans contact.

\paragraph{Avantages :}
\begin{itemize}
    \item pas d'usure mécanique ;
    \item cadence élevée ;
    \item bonne répétabilité.
\end{itemize}

\subsubsection{Détecteur de proximité capacitif}

Il détecte la présence d'objets conducteurs ou isolants (plastique, verre, liquide).

\subsubsection{Détecteur photoélectrique}

Il utilise un faisceau lumineux pour détecter un objet.

Trois modes existent :

\begin{itemize}
    \item \textbf{barrage} : émetteur et récepteur séparés ;
    \item \textbf{réflex} : émetteur et récepteur dans le même boîtier, avec réflecteur ;
    \item \textbf{proximité} : l'objet renvoie lui-même le faisceau.
\end{itemize}

\subsubsection{Capteurs de grandeur physique}

Ces capteurs délivrent un signal analogique ou numérique.

\paragraph{Exemples :}
\begin{itemize}
    \item capteur de niveau (analogique ou TOR) ;
    \item capteur de température (sonde PT100, thermocouple) ;
    \item capteur de pression ;
    \item codeur incrémental ou absolu.
\end{itemize}

\subsection{Caractéristiques d'un organe d'entrée}

Pour choisir un organe d'entrée, on tient compte de :

\begin{itemize}
    \item la \textbf{nature de la grandeur} détectée (position, niveau, température) ;
    \item le \textbf{type de signal} délivré (TOR, analogique, numérique) ;
    \item la \textbf{tension d'alimentation} (24 V DC, 230 V AC) ;
    \item la \textbf{distance de détection} ;
    \item la \textbf{cadence} de fonctionnement ;
    \item l'\textbf{environnement} (poussière, humidité, température) ;
    \item le \textbf{type de sortie} (PNP, NPN, contact sec).
\end{itemize}

\subsection{Raccordement des organes d'entrée à l'API}

Les organes d'entrée sont raccordés aux \textbf{modules d'entrées} de l'API.

Chaque entrée possède :

\begin{itemize}
    \item une \textbf{adresse} (par exemple I0.0, I0.1, I0.2) ;
    \item un \textbf{type} (TOR, analogique) ;
    \item un \textbf{état logique} : 0 ou 1.
\end{itemize}

\paragraph{Exemple :}

\begin{table}[h]
\centering
\begin{tabularx}{\textwidth}{|c|X|c|}
\hline
\textbf{Adresse} & \textbf{Organe} & \textbf{Type} \\
\hline
I0.0 & Bouton Marche & TOR \\
\hline
I0.1 & Bouton Arrêt & TOR \\
\hline
I0.2 & Capteur de niveau bas & TOR \\
\hline
I0.3 & Capteur de niveau haut & TOR \\
\hline
I0.4 & Fin de course vérin & TOR \\
\hline
\end{tabularx}
\caption{Exemple de tableau des entrées}
\end{table}

\subsection{Exemple d'application : remplissage de cuve}

Dans une cuve automatique :

\begin{itemize}
    \item \textbf{Bouton Marche} : entrée I0.0 ;
    \item \textbf{Bouton Arrêt} : entrée I0.1 ;
    \item \textbf{Capteur niveau bas} : entrée I0.2 ;
    \item \textbf{Capteur niveau haut} : entrée I0.3.
\end{itemize}

L'API lit l'état de ces entrées à chaque cycle de scrutation et décide de l'ouverture ou de la fermeture de l'électrovanne.

\subsection{Erreurs fréquentes à éviter}

\begin{itemize}
    \item Confondre contact NO et NC ;
    \item Alimenter un capteur PNP comme un NPN ;
    \item Dépasser la tension maximale d'entrée de l'API ;
    \item Négliger la protection contre les surtensions ;
    \item Mal repérer les adresses des entrées ;
    \item Oublier de tester un capteur avant son raccordement.
\end{itemize}

\subsection{Conclusion}

Les organes d'entrée sont indispensables au fonctionnement d'un système automatisé. Ils informent la partie commande de l'état du système et de son environnement.

Leur choix dépend de la grandeur à détecter, du type de signal et des conditions d'utilisation. Un bon câblage et un repérage rigoureux sont essentiels pour éviter les erreurs de diagnostic.

\subsubsection*{À retenir}

\begin{quote}
Les organes d'entrée transforment une information physique en un signal exploitable par l'API. Ils constituent le premier maillon de la chaîne d'information.
\end{quote}

\paragraph{Schéma fondamental :}

\[
\text{Organes d'entrée} \rightarrow \text{API} \rightarrow \text{Organes de sortie}
\]

\[
\text{Capteurs, boutons, détecteurs} \rightarrow \text{Entrées API}
\]
% =========================================================
\section{Boutons-poussoirs et commutateurs}

\subsection{Introduction}

Dans un système automatisé, l'opérateur doit pouvoir intervenir pour :

\begin{itemize}
    \item mettre l'installation en marche ;
    \item arrêter la machine ;
    \item choisir un mode de fonctionnement ;
    \item signaler un défaut ou déclencher un arrêt d'urgence.
\end{itemize}

Les \textbf{boutons-poussoirs} et les \textbf{commutateurs} sont les organes de commande manuelle les plus utilisés dans les pupitres et armoires électriques.

Ils appartiennent à la \textbf{chaîne d'information} et sont raccordés aux entrées de l'API.

\subsection{Le bouton-poussoir}

\subsubsection{Définition}

Un bouton-poussoir est un organe de commande manuelle qui modifie l'état d'un contact lorsqu'on appuie dessus.

Il revient généralement à sa position initiale dès que l'opérateur relâche le bouton.

\subsubsection{Symbole}

Le symbole normalisé d'un bouton-poussoir est le suivant :

\begin{itemize}
    \item bouton à contact \textbf{NO} : contact ouvert au repos, fermé lorsqu'on appuie ;
    \item bouton à contact \textbf{NC} : contact fermé au repos, ouvert lorsqu'on appuie.
\end{itemize}

\subsubsection{Types de boutons-poussoirs}

On distingue :

\begin{itemize}
    \item le bouton \textbf{à contact NO} (Normalement Ouvert) ;
    \item le bouton \textbf{à contact NC} (Normalement Fermé) ;
    \item le bouton \textbf{à contacts NO + NC} (double contact) ;
    \item le bouton \textbf{à verrouillage} (reste enfoncé) ;
    \item le bouton \textbf{à clé} (commande protégée).
\end{itemize}

\subsection{Bouton à contact NO}

\subsubsection{Principe}

Un bouton à contact NO est ouvert au repos. Lorsqu'on appuie dessus, le contact se ferme et laisse passer le signal.

\[
\text{Repos} \rightarrow \text{Contact ouvert}
\qquad
\text{Appui} \rightarrow \text{Contact fermé}
\]

\subsubsection{Utilisation}

Il est utilisé pour les commandes de type \textbf{Marche}, \textbf{Départ cycle} ou \textbf{Acquittement}.

\paragraph{Exemple :} bouton « Marche » sur un pupitre de commande.

\subsection{Bouton à contact NC}

\subsubsection{Principe}

Un bouton à contact NC est fermé au repos. Lorsqu'on appuie dessus, le contact s'ouvre et interrompt le signal.

\[
\text{Repos} \rightarrow \text{Contact fermé}
\qquad
\text{Appui} \rightarrow \text{Contact ouvert}
\]

\subsubsection{Utilisation}

Il est utilisé pour les commandes de type \textbf{Arrêt} ou \textbf{Arrêt d'urgence}.

\paragraph{Exemple :} bouton « Arrêt » sur un pupitre de commande.

\paragraph{Important :} en cas de rupture de fil, un contact NC ouvre le circuit et provoque l'arrêt. C'est un principe de sécurité fondamental.

\subsection{Bouton d'arrêt d'urgence}

\subsubsection{Définition}

Le bouton d'arrêt d'urgence est un bouton à contact NC, généralement à \textbf{verrouillage mécanique}. Il reste enfoncé après l'appui et doit être déverrouillé manuellement (par rotation ou traction) avant de pouvoir redémarrer.

\subsubsection{Couleur}

Selon la norme, il est :

\begin{itemize}
    \item de couleur \textbf{rouge} ;
    \item sur fond \textbf{jaune}.
\end{itemize}

\subsubsection{Rôle}

Il a pour fonction d'arrêter immédiatement et en toute sécurité l'installation en cas de danger.

\paragraph{Attention :} l'arrêt d'urgence doit agir directement sur la chaîne d'énergie ou via un relais de sécurité, et non uniquement par le programme de l'API.

\subsection{Les commutateurs et sélecteurs}

\subsubsection{Définition}

Un commutateur ou sélecteur possède plusieurs positions \textbf{stables}. Il conserve son état après l'action de l'opérateur, contrairement au bouton-poussoir qui revient à sa position initiale.

\subsubsection{Exemples d'utilisation}

\begin{itemize}
    \item sélecteur « Manuel / Automatique » ;
    \item sélecteur « Marche / Arrêt » ;
    \item sélecteur de vitesse ;
    \item commutateur de mode de cycle.
\end{itemize}

\paragraph{Exemple :} un sélecteur à deux positions qui reste en position « Automatique » jusqu'à ce que l'opérateur le tourne vers « Manuel ».

\subsection{Différence entre bouton-poussoir et commutateur}

\begin{table}[h]
\centering
\begin{tabularx}{\textwidth}{|X|X|}
\hline
\textbf{Bouton-poussoir} & \textbf{Commutateur / sélecteur} \\
\hline
Retour automatique au repos & Conserve sa position \\
\hline
Action momentanée & Action permanente \\
\hline
Utilisé pour Marche, Arrêt, AU & Utilisé pour choix de mode \\
\hline
Un ou deux contacts & Plusieurs positions possibles \\
\hline
\end{tabularx}
\caption{Comparaison bouton-poussoir / commutateur}
\end{table}

\subsection{Contacts NO et NC : rappel}

\begin{itemize}
    \item \textbf{NO} : ouvert au repos, fermé lorsqu'on actionne ;
    \item \textbf{NC} : fermé au repos, ouvert lorsqu'on actionne.
\end{itemize}

\paragraph{Règle importante :} pour la sécurité, on utilise toujours des contacts \textbf{NC} pour les arrêts et les arrêts d'urgence.

\subsection{Raccordement à l'API}

Les boutons et commutateurs sont raccordés aux entrées \textbf{TOR} de l'API.

\paragraph{Exemple :}

\begin{table}[h]
\centering
\begin{tabularx}{\textwidth}{|c|X|c|}
\hline
\textbf{Adresse} & \textbf{Organe} & \textbf{Type de contact} \\
\hline
I0.0 & Bouton Marche & NO \\
\hline
I0.1 & Bouton Arrêt & NC \\
\hline
I0.2 & Bouton Arrêt d'urgence & NC verrouillé \\
\hline
I0.3 & Sélecteur Manuel/Auto & NO \\
\hline
I0.4 & Sélecteur Auto & NO \\
\hline
\end{tabularx}
\caption{Exemple de raccordement de boutons et sélecteurs}
\end{table}

\subsection{Exemple d'application : démarrage d'un moteur}

On souhaite démarrer un moteur à l'aide de deux boutons :

\begin{itemize}
    \item un bouton \textbf{Marche} (NO) ;
    \item un bouton \textbf{Arrêt} (NC).
\end{itemize}

Le programme de l'API doit :

\begin{itemize}
    \item démarrer le moteur lorsque le bouton Marche est appuyé ;
    \item maintenir le moteur en marche après relâchement (auto-maintien logiciel) ;
    \item arrêter le moteur lorsque le bouton Arrêt est appuyé.
\end{itemize}

\[
\text{Marche (NO)} \rightarrow \text{API} \rightarrow \text{Moteur}
\]

\[
\text{Arrêt (NC)} \rightarrow \text{API} \rightarrow \text{Arrêt du moteur}
\]

\subsection{Bonnes pratiques}

\begin{itemize}
    \item Repérer clairement chaque bouton sur le pupitre ;
    \item Utiliser des couleurs normalisées (vert pour Marche, rouge pour Arrêt) ;
    \item Utiliser un contact NC pour tout arrêt ou arrêt d'urgence ;
    \item Vérifier le fonctionnement du bouton avant le raccordement ;
    \item Noter soigneusement les adresses dans le tableau des entrées ;
    \item Éviter de cumuler plusieurs fonctions sur un même bouton.
\end{itemize}

\subsection{Erreurs fréquentes à éviter}

\begin{itemize}
    \item Confondre NO et NC ;
    \item Utiliser un bouton NO pour un arrêt d'urgence ;
    \item Raccorder un bouton sur une entrée déjà utilisée ;
    \item Oublier de respecter la tension d'alimentation du bouton ;
    \item Ne pas tester le bouton après câblage ;
    \item Ne pas distinguer bouton momentané et bouton verrouillé.
\end{itemize}

\subsection{Conclusion}

Les boutons-poussoirs et les commutateurs sont les organes de commande manuelle de base dans un système automatisé.

Ils permettent à l'opérateur d'agir sur la partie commande : mise en marche, arrêt, choix de mode, acquittement, arrêt d'urgence.

Leur raccordement correct aux entrées de l'API, ainsi que le respect des règles de sécurité, sont essentiels pour la fiabilité et la sécurité de l'installation.

\subsubsection*{À retenir}

\begin{quote}
Un bouton-poussoir est un organe momentané, tandis qu'un commutateur conserve sa position. Les contacts NO sont utilisés pour les commandes, les contacts NC pour les arrêts et les sécurités.
\end{quote}

\paragraph{Schéma fondamental :}

\[
\text{Boutons / commutateurs} \rightarrow \text{Entrées API} \rightarrow \text{Programme}
\]
% =========================================================
\section{Capteurs, détecteurs et fins de course}

\subsection{Introduction}

Dans un système automatisé, la partie commande doit connaître à tout moment l'état de la partie opérative et de son environnement.

Pour cela, elle utilise des \textbf{capteurs} et des \textbf{détecteurs}, qui traduisent une information physique en un signal exploitable par l'API.

Parmi ces organes, les \textbf{fins de course} occupent une place importante : ils indiquent la position extrême d'un organe mobile.

\subsection{Définition d'un capteur}

Un capteur est un composant qui prélève une information sur une grandeur physique (position, présence, niveau, température, pression, etc.) et la transforme en un signal électrique exploitable par la partie commande.

\[
\text{Grandeur physique} \rightarrow \text{Capteur} \rightarrow \text{Signal électrique} \rightarrow \text{API}
\]

\subsection{Classification des capteurs}

On peut classer les capteurs selon plusieurs critères.

\subsubsection{Selon le type de signal}

\begin{itemize}
    \item \textbf{Capteurs TOR} (Tout Ou Rien) : délivrent un signal 0 ou 1 ;
    \item \textbf{Capteurs analogiques} : délivrent un signal proportionnel à la grandeur mesurée (0--10 V, 4--20 mA) ;
    \item \textbf{Capteurs numériques} : délivrent un code (binaire, Gray, BCD).
\end{itemize}

\subsubsection{selon le mode de détection}

\begin{itemize}
    \item \textbf{Avec contact} : fin de course mécanique ;
    \item \textbf{Sans contact} : détecteurs de proximité, photoélectriques, ultrasoniques.
\end{itemize}

\subsubsection{Selon la nature de la grandeur}

\begin{itemize}
    \item Détection de présence ou de position ;
    \item Mesure de niveau ;
    \item Mesure de température ;
    \item Mesure de pression ;
    \item Mesure de débit ;
    \item Mesure de vitesse ou de position angulaire.
\end{itemize}

\subsection{Les fins de course}

\subsubsection{Définition}

Un fin de course est un capteur mécanique qui détecte la position extrême d'un organe mobile (vérin, porte, tiroir, etc.).

Il est généralement équipé d'un contact \textbf{NO} ou \textbf{NC} et d'un levier ou d'un galet actionné par l'organe en mouvement.

\subsubsection{Principe de fonctionnement}

Lorsque l'organe mobile arrive en fin de course, il appuie sur le levier du fin de course, ce qui ouvre ou ferme le contact.

\[
\text{Organe mobile en butée} \rightarrow \text{Action sur le levier} \rightarrow \text{Contact ouvert ou fermé}
\]

\subsubsection{Types de fins de course}

\begin{itemize}
    \item Fin de course \textbf{à levier} ;
    \item Fin de course \textbf{à galet} ;
    \item Fin de course \textbf{à tige} ;
    \item Fin de course \textbf{à billes}.
\end{itemize}

\subsubsection{Utilisation}

Les fins de course servent à :

\begin{itemize}
    \item détecter la position d'un vérin (sorti/rentré) ;
    \item limiter la course d'un organe mobile ;
    \item arrêter un moteur en fin de course ;
    \item confirmer une position avant d'autoriser l'étape suivante d'un cycle.
\end{itemize}

\paragraph{Exemple :} dans un vérin pneumatique, deux fins de course indiquent respectivement la position \textbf{rentrée} et la position \textbf{sortie}.

\subsection{Les détecteurs de proximité}

\subsubsection{Définition}

Un détecteur de proximité est un capteur \textbf{sans contact} qui détecte la présence d'un objet à faible distance.

\subsubsection{Détecteur inductif}

Il détecte uniquement les objets \textbf{métalliques}.

\paragraph{Principe :} il crée un champ magnétique alternatif et détecte la perturbation provoquée par un objet métallique qui s'approche.

\paragraph{Avantages :}
\begin{itemize}
    \item pas d'usure mécanique ;
    \item cadence élevée ;
    \item insensible à la poussière et à l'humidité.
\end{itemize}

\paragraph{Utilisation :} détection de cames, de pignons, de pièces métalliques.

\subsubsection{Détecteur capacitif}

Il détecte les objets \textbf{conducteurs ou isolants} (plastique, verre, bois, liquides).

\paragraph{Principe :} il détecte la variation de capacité provoquée par l'approche d'un objet.

\paragraph{Utilisation :} détection de niveau de liquide à travers une paroi, détection de pièces plastiques.

\subsubsection{Détecteur photoélectrique}

Il utilise un \textbf{faisceau lumineux} (infrarouge ou visible) pour détecter un objet.

Trois modes existent :

\begin{itemize}
    \item \textbf{Barrage} : émetteur et récepteur séparés ; l'objet coupe le faisceau.
    \item \textbf{Réflex} : émetteur et récepteur dans le même boîtier, avec un réflecteur ; l'objet coupe le faisceau.
    \item \textbf{Proximité} : l'objet lui-même renvoie le faisceau vers le récepteur.
\end{itemize}

\paragraph{Utilisation :} détection de pièces, comptage, sécurité (barrière immatérielle).

\subsubsection{Détecteur à ultrasons}

Il émet des ondes ultrasonores et mesure le temps de retour de l'écho.

\paragraph{Utilisation :} détection d'objets transparents, mesure de distance, détection de niveau.

\subsection{Les capteurs de grandeur physique}

\subsubsection{Capteurs de niveau}

Ils détectent ou mesurent le niveau d'un liquide ou d'un solide dans un réservoir.

\begin{itemize}
    \item \textbf{Niveau TOR} : flotteur, électrode, détecteur capacitif ;
    \item \textbf{Niveau analogique} : capteur de pression, ultrasons, radar.
\end{itemize}

\subsubsection{Capteurs de température}

\begin{itemize}
    \item \textbf{Thermocouple} : mesure de hautes températures ;
    \item \textbf{Sonde PT100} : mesure précise et linéaire ;
    \item \textbf{Thermistance} : mesure économique.
\end{itemize}

\subsubsection{Capteurs de pression}

Ils mesurent la pression d'un fluide (air, eau, huile).

\subsubsection{Codeurs}

Un codeur est un capteur qui mesure une position angulaire ou une vitesse de rotation.

\begin{itemize}
    \item \textbf{Codeur incrémental} : délivre des impulsions ; il faut un comptage ;
    \item \textbf{Codeur absolu} : délivre directement la position sous forme de code.
\end{itemize}

\subsection{Caractéristiques d'un capteur}

Pour choisir un capteur, on tient compte de :

\begin{itemize}
    \item la \textbf{grandeur} à détecter ;
    \item le \textbf{type de signal} de sortie (TOR, analogique, numérique) ;
    \item la \textbf{distance de détection} ;
    \item la \textbf{cadence} de fonctionnement ;
    \item la \textbf{tension d'alimentation} (24 V DC, 230 V AC) ;
    \item le \textbf{type de sortie} (PNP, NPN, contact sec) ;
    \item l'\textbf{environnement} (poussière, humidité, température, vibrations) ;
    \item l'\textbf{indice de protection} (IP).
\end{itemize}

\subsection{Sorties PNP et NPN}

Les détecteurs de proximité électroniques possèdent souvent une sortie à transistor.

\begin{itemize}
    \item \textbf{Sortie PNP} : le transistor relie la sortie au \textbf{+24 V} lorsqu'il est actif ;
    \item \textbf{Sortie NPN} : le transistor relie la sortie au \textbf{0 V} lorsqu'il est actif.
\end{itemize}

\paragraph{Attention :} il faut respecter le type de sortie attendu par l'entrée de l'API. Un capteur PNP ne se raccorde pas comme un capteur NPN.

\subsection{Raccordement à l'API}

Les capteurs et détecteurs sont raccordés aux entrées de l'API.

\paragraph{Exemple de tableau des entrées :}

\begin{table}[h]
\centering
\begin{tabularx}{\textwidth}{|c|X|c|c|}
\hline
\textbf{Adresse} & \textbf{Organe} & \textbf{Type} & \textbf{Sortie} \\
\hline
I0.0 & Fin de course vérin rentré & TOR & Contact NC \\
\hline
I0.1 & Fin de course vérin sorti & TOR & Contact NO \\
\hline
I0.2 & Détecteur inductif pièce & TOR & PNP \\
\hline
I0.3 & Détecteur photoélectrique & TOR & PNP \\
\hline
I0.4 & Capteur de niveau bas & TOR & Contact NO \\
\hline
IW64 & Capteur de pression & Analogique & 4--20 mA \\
\hline
\end{tabularx}
\caption{Exemple de raccordement de capteurs}
\end{table}

\subsection{Exemple d'application : détection de position d'un vérin}

Un vérin pneumatique doit être commandé en va-et-vient.

Deux capteurs de position sont installés :

\begin{itemize}
    \item un capteur \textbf{vérin rentré} (fin de course) ;
    \item un capteur \textbf{vérin sorti} (fin de course).
\end{itemize}

Le programme de l'API utilise ces informations pour :

\begin{itemize}
    \item autoriser la sortie du vérin uniquement lorsqu'il est rentré ;
    \item autoriser la rentrée uniquement lorsqu'il est sorti.
\end{itemize}

\[
\text{Vérin rentré} \rightarrow \text{API} \rightarrow \text{Commande sortie}
\]

\[
\text{Vérin sorti} \rightarrow \text{API} \rightarrow \text{Commande rentrée}
\]

\subsection{Exemple d'application : comptage de pièces}

Un détecteur photoélectrique détecte le passage de pièces sur un convoyeur.

Chaque détection est comptée par l'API. Lorsque le compteur atteint une valeur déterminée, l'API arrête le convoyeur ou déclenche une alarme.

\[
\text{Impulsion détecteur} \rightarrow \text{Compteur API} \rightarrow \text{Action}
\]

\subsection{Bonnes pratiques}

\begin{itemize}
    \item Vérifier la compatibilité PNP/NPN avec l'API ;
    \item Respecter la tension d'alimentation du capteur ;
    \item Éviter les environnements trop perturbés (vibrations, poussière) ;
    \item Repérer clairement chaque capteur sur le schéma ;
    \item Utiliser un contact NC pour les sécurités ;
    \item Tester chaque capteur avant raccordement ;
    \item Protéger les câbles contre les perturbations électromagnétiques.
\end{itemize}

\subsection{Erreurs fréquentes à éviter}

\begin{itemize}
    \item Confondre capteur PNP et NPN ;
    \item Utiliser un capteur inductif pour détecter du plastique ;
    \item Dépasser la distance de détection nominale ;
    \item Négliger l'indice de protection dans un environnement humide ;
    \item Mal régler la sensibilité d'un détecteur photoélectrique ;
    \item Ne pas tenir compte du temps de réponse du capteur.
\end{itemize}

\subsection{Conclusion}

Les capteurs, détecteurs et fins de course sont les organes qui informent la partie commande sur l'état du système.

Leur choix dépend de la grandeur à détecter, du type de signal, de l'environnement et du type de sortie attendu par l'API.

Un bon raccordement et un repérage rigoureux sont essentiels pour la fiabilité du système automatisé.

\subsubsection*{À retenir}

\begin{quote}
Un capteur transforme une grandeur physique en signal exploitable par l'API. Un fin de course détecte la position extrême d'un organe mobile. Un détecteur de proximité fonctionne sans contact.
\end{quote}

\paragraph{Schéma fondamental :}

\[
\text{Grandeur physique} \rightarrow \text{Capteur} \rightarrow \text{Signal} \rightarrow \text{Entrée API}
\]
\end{document}